\section{重点问题解决记录}

\begin{longtable}[htbp!]{ m{3.5cm}<{\centering}  m{3.5cm}<{\centering}m{4.5cm}<{\centering}m{1cm}<{\centering} m{1.5cm}<{\centering}}

\hline

问题描述    &   问题产生原因    &   问题解决方案\&实际解决效果  &   机器人版本号    &   解决人员   \\
\hline
\endfirsthead
\hline

问题描述    &   问题产生原因    &   问题解决方案\&实际解决效果  &   机器人版本号    &   解决人员   \\
\hline
\endhead
\hline
\multicolumn{2}{r}{接下页}
\endfoot
\caption{重点问题解决记录}
\label{tab:重点问题解决记录}
\endlastfoot
%   填写这里
外部弹仓在静止状态下也存在大量弹丸无法下落问题 & 外部弹仓底板倾斜角度不足 & 调大外部弹仓底板倾斜角度，落弹率提高至可接受范围 & v0.8 & (机械) XXX \\
\hline
三位一体固定座与Yaw轴的连接刚性不足 & 连接方式采用$\qty{4}{\mm}$玻纤板，抗弯刚性不足 & 将$\qty{4}{\mm}$玻纤板更换为$\qty{8}{\mm}$碳管，同时连接处的各铝柱用打印件嵌套，整体刚性提高至可接受范围 & v1.2 & (机械)XXX \\
\hline
底盘减速箱频繁损坏 & 外包的铝件质量出现问题 & 更换全新的铝件，后续不再出现损坏现象 & 1.4 & (机械)XXX \\
\hline
弹仓容弹量增大后仅有约550发，仍然无法满足750发容弹量的需求 & 22赛季步兵设计紧凑，仅增大弹仓至极限尺寸也无法达到需求 & 确定轮系不做修改的前提后，扩大了云台安装座固定铝管的间距与保护框的大小，使得弹仓可修改的空间更大。目前极限可容纳1150颗弹丸，完全满足需求 & v1.6 & (机械)XXX \\
\hline
TPU弹丸定心方案发射测试中，重复度效果不理想，双发问题较为严重 & 结合后续的测试结论得出，关键因素应是弹链设计部分未充分考虑Pitch轴活动关节处弹丸阻力的影响，弹丸在接触到单发限位时，Pitch轴活动关节处会对整条弹链产生的阻力极大，而非TPU的问题 & 结合整体进度考虑，我们尝试通过减小阻力以减缓双发问题的发生频率，通过更换较成熟且开源资料更多的双轴承定心方案以及拉簧-打印件弹丸定心方案去提高重复度 & v2.0 & (机械)XXX \\
\hline
双轴承弹丸定心、拉簧-打印件弹丸定心方案发射测试中，云台存在频繁双发且低头卡弹的现象 & 弹链阻力较大，中心供弹推弹力度不足，单发限位阻力较大，链路设计可优化 & 在上滑块内侧以及外部弹链固定板接近Pitch活动关节处添加轴承，减少阻力。同时增大中心供弹齿轮传动减速比。减小双轴承间压缩量(双轴承弹丸定心方案),减小拉簧K值(拉簧-打印件定心方案)以减小单发限位的阻力，目前双发情况基本解决，重复度满足需求 & v2.8 & (机械)XXX \\
\hline
在视觉标签的法线向量与相机光心向量之间所成角度太大时，我们的算法会出现提取标签内容发生变形的问题 & 因为在进行提取内容时需要涉及取模操作 & 在提取标签区域时计算旋转角度再进行仿射变换可以很好的规避这个问题 & v2.0 & (视觉)XXX \\
\hline
NUC网口在接入MID-360数据线时，会导致无法通过有线连接的方式接入NUC & MID-360会占用有线网卡  & 使用3口交换机连接NUC，同时接MID-360，此时可以通过网线调试机器人  & v 2.0 & (控制)XXX \\
\hline


\end{longtable}